\section{Introduction}

Consider the stochastic optimization problem
\[
\min_{\theta\in\mathbb R^p}F(\theta),\qquad
F(\theta)=\mathbb E\ell(\theta;Z),
\]
when the latent input $Z$ is observed only through $W=Z+U$.
Even for centered measurement error, the score
$g(\theta;W)=\nabla_\theta\ell(\theta;W)$ is generally biased for
$\nabla F(\theta)$. A decreasing step size controls stochastic
fluctuations but does not eliminate this discrepancy. The question is
whether noisy observations can nevertheless support an explicit
stochastic-gradient method for the clean objective.

For Gaussian error with known covariance, the classical complex
corrected-score identity provides a construction
\cite{stefanski1989,novick2002}. If the score has an integrable entire
extension, an independent imaginary Gaussian perturbation cancels the
measurement error in expectation. This leads to complex-corrected
stochastic gradient descent (CCSGD):
\[
\theta_{t+1}=\theta_t-\eta_t\operatorname{Re}
g^{\mathbb C}(\theta_t;W_t+iV_t),\qquad
V_t\sim N(0,\Omega).
\]
Each iteration uses one noisy observation and one complex score evaluation.
Only the noise covariance is required; it may be correlated or singular.

The central analytical issue is control of the corrected gradient along
the stochastic trajectory. Pointwise unbiasedness and finite variance
do not by themselves give the conditional moment bounds required for
convergence. We address this issue through three results.
\begin{enumerate}
\item \emph{Exact correction.}
We formulate the Gaussian correction for vector-valued entire scores
under an explicit exponential growth condition. The construction is
independent of the convexity of the objective.
\item \emph{Relative variance control.}
We bound the correction error using the score's input derivative and
the tails of the latent and Gaussian variables. The derivative amplitude
may grow with the parameter, provided its square is controlled by the
squared objective gradient. Conditional orthogonality then yields the
classical SGD relative variance condition and local higher-moment control
for limit laws.
\item \emph{Convergence and fluctuations.}
For polynomially decaying step sizes with exponent $1/2<\alpha\le1$,
we establish convex objective bounds, strongly convex mean-square rates
and almost-sure convergence, and nonconvex stationarity bounds.
Under local regularity, we identify the last-iterate and averaged
Gaussian limits, including the $\alpha=1$ endpoint.
\end{enumerate}
The analysis connects the classical Gaussian identity to stochastic
approximation through explicit moment estimates
\cite{novick2002,fabian1968,polyakjuditsky1992}. Examples and numerical
studies illustrate the method; a separate moment argument covers
nonlinear-feature regression.
